% ============================================================================
% Fichier  : partie2/P02_C06_chartes_rgpd.tex
% Titre    : Chartes, Conformité et Protection des Données
% Auteur   : MINKA MI NGUIDJOÏ Thierry Emmanuel
% ============================================================================
% Sources : Loi camerounaise 2024/017 ; Convention de Malabo 2014 ;
% ISO/IEC 27701. Mairie d'Eséka : cas fictif.
% ============================================================================

\chapter{Chartes, Conformité et Protection des Données}
\label{chap:conformite}
\chaptermark{Chartes et conformité}

\epigraph{%
    « Les données personnelles ne sont pas la propriété de la mairie~:\\
    elles lui sont confiées. »%
}{Principe du cours}

\niveauChapitre{Fondamental}{Le Chapitre~\ref{chap:politique}.}

\begin{boxobjectifs}
\begin{itemize}
    \item Rédiger et faire vivre une charte informatique.
    \item Énoncer les principes de protection des données personnelles.
    \item Situer la loi camerounaise 2024/017 et la Convention de Malabo.
    \item Définir le rôle du délégué à la protection des données (\sigle{DPO}).
    \item Distinguer sécurité et conformité, et montrer leur recouvrement.
\end{itemize}
\end{boxobjectifs}

\bigskip

% ============================================================================
\section*{Question directrice : la mairie a-t-elle le droit de faire ce qu'elle fait de ses données ?}
% ============================================================================

La mairie détient des données très sensibles~: naissances, mariages, décès,
contribuables. En détenir ne vaut pas en faire ce que l'on veut. Ce chapitre
pose le cadre de conformité et l'engagement des agents.

\medskip
\textbf{Plan du chapitre.}
\begin{enumerate}
    \item La charte informatique (section~\ref{sec:c06-charte}).
    \item Les principes de protection des données (section~\ref{sec:c06-principes}).
    \item Le cadre juridique camerounais (section~\ref{sec:c06-droit}).
    \item Le délégué à la protection des données (section~\ref{sec:c06-dpo}).
    \item Sécurité et conformité (section~\ref{sec:c06-recouvrement}).
\end{enumerate}

\separateur

% ============================================================================
\section{La Charte Informatique}
\label{sec:c06-charte}
% ============================================================================

\begin{boxdef}
\termecle{Charte informatique}~: document opérationnel qui précise les règles
d'utilisation du système d'information, les droits et devoirs des utilisateurs,
et les sanctions en cas de non-respect. Signée, elle engage l'agent.
\end{boxdef}

La charte traduit la politique en obligations individuelles~: usage des postes,
gestion des mots de passe, messagerie, supports amovibles, télétravail. Sa
force vient de la signature~: un agent qui a signé ne peut plaider
l'ignorance.

% ============================================================================
\section{Les Principes de Protection des Données}
\label{sec:c06-principes}
% ============================================================================

\begin{tableaucours}{L{3.6cm}L{8.6cm}}{%
    \tableauHeader{Principe} &
    \tableauHeader{Exigence}
}
Licéité et loyauté & Une base légale pour chaque traitement \\
Finalité & Données collectées pour un but déterminé \\
Minimisation & Données pertinentes et limitées au nécessaire \\
Exactitude & Données exactes et tenues à jour \\
Conservation limitée & Durée de conservation appropriée \\
Sécurité & Mesures techniques et organisationnelles \\
Responsabilité & Démontrer la conformité \\
\end{tableaucours}
\vspace{-0.8em}
\begin{center}{\small\itshape Tableau~6.1. Les principes de protection des données}\end{center}

\begin{boxobjet}{Le fichier des contribuables}
\textbf{Constatation.} Le fichier des contribuables contient des informations
qui ne servent plus (anciens redevables partis depuis dix ans).

\medskip
\textbf{Analyse.} Deux principes sont en défaut~: la \emph{minimisation} et la
\emph{conservation limitée}. Conserver sans fin expose inutilement des données
personnelles. La conformité impose une politique de purge.
\end{boxobjet}

% ============================================================================
\section{Le Cadre Juridique Camerounais}
\label{sec:c06-droit}
% ============================================================================

\begin{boxjuris}
La loi camerounaise n°~2024/017 du 23~décembre~2024 régit la protection des
données à caractère personnel~: elle impose une base légale, la sécurité des
données et des droits aux personnes concernées. La Convention de l'Union
africaine sur la cybersécurité et la protection des données (Convention de
\textbf{Malabo}, 2014) fixe le cadre continental. La mairie, responsable de
traitement, doit s'y conformer.
\end{boxjuris}

La conformité n'est pas optionnelle pour une administration~: elle conditionne
la légalité même du service rendu.

% ============================================================================
\section{Le Délégué à la Protection des Données}
\label{sec:c06-dpo}
% ============================================================================

\begin{boxdef}
\termecle{Délégué à la protection des données} (\sigle{DPO})~: personne chargée
de veiller à la conformité des traitements de données personnelles, de
conseiller l'organisation et de faire le lien avec l'autorité de protection.
\end{boxdef}

À Eséka, faute de \sigle{DPO} interne, la mairie confie ce rôle à
Me~\textsc{Essomba}, conseil extérieur. Le \sigle{DPO} tient le registre des
traitements, évalue les risques pour la vie privée et instruit les demandes des
usagers.

\begin{boxmetier}{Responsable conformité}
\textbf{Conformité n'est pas sécurité.} Le responsable conformité vérifie que
la mairie respecte ses obligations~; le \sigle{RSSI} réduit les risques. Les
deux se recoupent mais ne se confondent pas~: on peut être sécurisé et non
conforme (données gardées trop longtemps), ou conforme sur le papier et peu
sécurisé.
\end{boxmetier}

% ============================================================================
\section{Sécurité et Conformité}
\label{sec:c06-recouvrement}
% ============================================================================

La sécurité (\sigle{DIC}) et la conformité (respect des règles) se recouvrent
largement~: la loi exige la sécurité, et la sécurité sert la conformité. Mais
elles diffèrent de logique~: la sécurité part du \emph{risque}, la conformité
part de l'\emph{obligation}. Un bon \sigle{SMSI} tient les deux.

\begin{boxinfo}
\textbf{Un pont vers le reste du cours.} La conformité donne des exigences
d'entrée à l'appréciation des risques (partie IV)~: certaines mesures seront
obligatoires quel que soit le niveau de risque, parce que la loi les impose.
\end{boxinfo}

\separateur

\begin{boxresume}
\textbf{Ce qu'il faut retenir de ce chapitre~:}
\begin{itemize}
    \item La \textbf{charte informatique} engage l'agent par sa signature.
    \item Les \textbf{principes} de protection des données (finalité,
          minimisation, conservation limitée\ldots) encadrent les traitements.
    \item La \textbf{loi 2024/017} et la \textbf{Convention de Malabo} fixent
          le cadre camerounais et africain.
    \item Le \textbf{DPO} veille à la conformité des traitements.
    \item \textbf{Sécurité et conformité} se recoupent sans se confondre.
\end{itemize}
\end{boxresume}

\begin{boxinfo}
\textbf{Points de contrôle de l'élève.} Je sais~: rédiger une règle de charte~;
citer les principes de protection des données~; situer la loi 2024/017 et
Malabo~; définir le rôle du \sigle{DPO}~; distinguer sécurité et conformité.
\end{boxinfo}

\bigskip

\begin{boxexo}[title={\ding{45}\quad Exercice~6.1. La charte \niveaubase}]
Citez trois règles que doit contenir la charte informatique de la mairie.
\end{boxexo}

\begin{boxexo}[title={\ding{45}\quad Exercice~6.2. Les principes \niveaubase}]
Énoncez quatre principes de protection des données personnelles.
\end{boxexo}

\begin{boxexo}[title={\ding{45}\quad Exercice~6.3. La loi \niveaubase}]
Quelle loi camerounaise régit la protection des données personnelles, et de
quand date-t-elle~?
\end{boxexo}

\begin{boxexo}[title={\ding{45}\quad Exercice~6.4. Deux principes en défaut \niveaumoyen}]
Pour le fichier des contribuables conservé sans fin, quels principes sont
violés~? Que faire~?
\end{boxexo}

\begin{boxexo}[title={\ding{45}\quad Exercice~6.5. Rôle du DPO \niveaumoyen}]
Décrivez trois missions du \sigle{DPO} à la mairie d'Eséka.
\end{boxexo}

\begin{boxexo}[title={\ding{45}\quad Exercice~6.6. Sécurisé mais non conforme \niveaumoyen}]
Donnez un exemple d'une mairie bien sécurisée mais non conforme à la loi.
\end{boxexo}

\begin{boxexo}[title={\ding{45}\quad Exercice~6.7. Base légale \niveaumoyen}]
Quelle base légale justifie que la mairie traite les actes d'état civil~?
\end{boxexo}

\begin{boxexo}[title={\ding{45}\quad Exercice~6.8. Signature \niveaumoyen}]
Pourquoi la signature de la charte est-elle juridiquement importante~?
\end{boxexo}

\begin{boxexo}[title={\ding{45}\quad Exercice~6.9. Charte rédigée \niveauexpert}]
Rédigez une charte informatique d'une page pour les agents de la mairie
d'Eséka.
\end{boxexo}

\begin{boxexo}[title={\ding{45}\quad Exercice~6.10. Registre des traitements \niveauexpert}]
Dressez un extrait du registre des traitements de la mairie (trois traitements,
finalité, données, durée de conservation).
\end{boxexo}

\bigskip

\begin{boxinfo}
\textbf{Pour aller plus loin.}
\begin{itemize}
    \item Loi 2024/017~\cite{loi2024017}~; Convention de Malabo~\cite{malabo2014}.
\end{itemize}
\end{boxinfo}

\begin{boxinfo}
\textbf{Transition.} Gouvernance, politique et conformité posées, le cours
entre dans la norme qui les structure~: le Chapitre~\ref{chap:iso27000}
présente la famille \norme{ISO/IEC 27000}.
\end{boxinfo}
